\chapter*{INTRODUCTION GÉNÉRALE}
\markboth{\MakeUppercase{INTRODUCTION GÉNÉRALE}}{}
\addcontentsline{toc}{chapter}{INTRODUCTION GÉNÉRALE}
\adjustmtc
\thispagestyle{MyStyle}
\begin{sloppypar} 
\vspace{-13mm}
Avec la croissance exponentielle du marché des applications mobiles, le développement d’applications dédiées à divers besoins est crucial. Dans le domaine de l’éducation, notamment pour les étudiants universitaires, l’accès à des informations pertinentes et la gestion efficace des services posent des défis significatifs.

\vspace{-1mm}
Actuellement, les étudiants font face à plusieurs défis majeurs pour gérer leur vie universitaire, notamment la difficulté à trouver des logements près de l'université et à accéder aux installations comme les bibliothèques et les cafés, souvent dispersées sur différentes plateformes. De plus, la recherche d'informations essentielles comme les événements, les stages, et les calendriers d'examens est complexe, nécessitant la consultation de multiples systèmes et applications.

\vspace{-1mm}
Notre solution consiste à développer une plateforme centralisée qui rassemble toutes ces données sous un même toit, permettant aux étudiants d'accéder facilement et rapidement à des informations fiables sur les logements disponibles et les facilités. En centralisant ces informations, nous visons à simplifier et optimiser significativement l'expérience des étudiants dans la gestion de leurs activités académiques et personnelles.

\vspace{-1mm}
Ce rapport détaille neuf chapitres : \textbf{Le premier} présente l’organisme d’accueil et le projet, avec étude de l’existant et solution proposée. \textbf{Le deuxième} traite de l'architecture logicielle, des outils utilisés et de la méthodologie de gestion de projet. \textbf{Le troisième} couvre les stratégies de test et le processus de déploiement. \textbf{Le quatrième} aborde la gestion des utilisateurs. \textbf{Le cinquième} explore la gestion des universités, des lieux et des examens. \textbf{Le sixième} examine la gestion des transports, des événements et des stages. \textbf{Le septième} traite des publications et des notifications. \textbf{Le huitième} explore les cartes, les conversations en ligne et les tableaux de bord. Une conclusion évalue le travail accompli et les futures perspectives.

\vspace{-1mm}
En conclusion, nous présenterons une synthèse de notre travail ainsi que les perspectives d'extension de ce projet.
\end{sloppypar}
\endinput
